% concatenated from N-pa-one.tex
\documentclass[11pt]{article}
\usepackage{amsmath, amsthm, amssymb, geometry, enumitem, graphicx, booktabs}
\usepackage{hyperref}
\usepackage{mathtools}
%%%%%%% adding nice affiliation command and orcid
\usepackage{authblk}
\usepackage{orcidlink}


\geometry{margin=1in}

\title{Path--Averaged Contractions: A New Generalization of the Banach Contraction Principle\footnote{MSC2020: 47H10, 54E50 \\Keywords: Path Averaged Contractions, contractive mapping, fixed point, complete metric space.}}
\author{Nicola Fabiano\, \orcidlink{0000-0003-1645-2071}}
%0000-0003-1645-2071
\affil{``Vin\v{c}a'' Institute of Nuclear Sciences - National 
Institute of the Republic of Serbia, University of Belgrade, Mike Petrovi\'{c}a 
Alasa 12--14, 11351 Belgrade, Serbia; nicola.fabiano@gmail.com}
\date{}
%%
%% 2020 Mathematics Subject Classification. Primary 47H10; Secondary 54E50.
%%Keywords. Path Averaged Contractions, contractive mapping, fixed point, 
%%complete metric space.
%%
%%


\newtheorem{theorem}{Theorem}
\newtheorem{lemma}[theorem]{Lemma}
\newtheorem{proposition}[theorem]{Proposition}
\newtheorem{definition}[theorem]{Definition}
\newtheorem{example}[theorem]{Example}
\newtheorem{remark}[theorem]{Remark}
\newtheorem{corollary}[theorem]{Corollary}


\begin{document}

\maketitle




%%%%%%%%%%%%%%%%%%%%%%%%%%%%%%%%%%%%%%%%%%%%%%%%%%%%%%%%%%%%%

%%%%%%%%%%%%%%%%%%%%%%%%%%%%%%%%%%%%%%%%%%%%%%%%%%

\begin{abstract}
We introduce a novel class of self-mappings on metric spaces, called \textbf{PA-contractions} (Path-Averaged Contractions), defined by an averaging condition over iterated distances. We prove that every continuous PA-contraction on a complete metric space has a unique fixed point, and the Picard iterates converge to it. This condition strictly generalizes the classical Banach contraction principle. We provide examples showing that PA-contractions are independent of F-contractions, Kannan, Chatterjea, and \'Ciri\'c contractions. A comparison table highlights the distinctions. The PA-condition captures long-term contractive behavior even when pointwise contraction fails.
\end{abstract}
\maketitle




\section{Introduction}

The Banach contraction principle \cite{Banach1922} is a cornerstone of nonlinear analysis. It states that if $ (X,d) $ is a complete metric space and $ T: X \to X $ satisfies
\begin{equation}
d(Tx, Ty) \leq k d(x,y), \quad \text{for some } k \in (0,1),
\end{equation}
then $ T $ has a unique fixed point. Numerous generalizations have been proposed, including Kannan \cite{Kannan1968}, \'Ciri\'c \cite{Ciric1971,Ciricorbit,Ciric1974,Ciric1979}, 
Chatterjea~\cite{Chatterjea1972},  with a systematic comparison of early variants provided by Rhoades \cite{Rhoades1977}, who analyzed over forty types of contractive mappings.
More recently, \textbf{F-contractions} introduced by Wardowski \cite{Wardowski2012}, extended and refined by several
authors~\cite{fabiano2022,Hussain2012,Piri2014,Secelean2013}.

In this paper, we propose a fundamentally different generalization: instead of pointwise or functional inequalities, we impose a \textit{mean contraction condition} over the orbit of pairs. This leads to the notion of \textbf{PA-contraction}, which depends on the cumulative behavior of iterates and is not reducible to known classes. For general background on fixed point theory, see the comprehensive text by Agarwal, Meehan, and O'Regan \cite{Agarwal2011}.

We prove a fixed point theorem for PA-contractions, show they generalize Banach contractions, and provide counterexamples demonstrating independence from other classes. Finally, we include a comparison table of major contraction types.

\section{PA-Contractions}

We begin with the central definition.

\begin{definition}[PA-Contraction]
Let $ (X,d) $ be a metric space. A mapping $ T: X \to X $ is called a \textbf{PA-contraction} (Path-Averaged Contraction) if there exists $ \alpha \in (0,1) $ 
and $ N \in \mathbb{N}$ such that for all $ x, y \in X $ and all $ n \ge N$,
\begin{equation}
\frac{1}{n} \sum_{k=0}^{n-1} d(T^{k+1}x, T^{k+1}y) \leq \alpha \cdot \frac{1}{n} \sum_{k=0}^{n-1} d(T^k x, T^k y).
\label{pa1}
\end{equation}
Equivalently, in sum form:
\begin{equation}
\sum_{k=0}^{n-1} d(T^{k+1}x, T^{k+1}y) \leq \alpha \sum_{k=0}^{n-1} d(T^k x, T^k y).
\label{pa2}
\end{equation}
\end{definition}

The condition requires that the average distance between successive iterates contracts by a uniform factor $ \alpha < 1 $, regardless of the starting pair $ (x,y) $.

The PA-contraction is loosely based on the path integral formulation of 
quantum mechanics~\cite{fabianopathintegral}, where the role of the path of the particle is taken by the orbit of the mapping $T$. 

\begin{remark}
If $ N = 1 $, then setting $ n = 1 $ in \eqref{pa1} yields $ d(Tx, Ty) \leq \alpha d(x, y) $, so $ T $ is a Banach contraction. However, for $ N > 1 $, the condition allows transient expansion, and the mapping need not satisfy any pointwise contraction. The class of PA-contractions thus \textbf{strictly generalizes} Banach contractions, as shown by Example~\ref{ex:panotbanach}.
The PA-condition is asymptotic in nature, depending on long-term behavior rather than immediate contraction.
\end{remark}

\section{Fixed Point Theorem}

Our main result is the following.

\begin{theorem}[Fixed Point Theorem for PA-Contractions]
Let $ (X,d) $ be a complete metric space, and let $ T: X \to X $ be a continuous PA-contraction. Then $ T $ has a unique fixed point $ x^* \in X $, and for any $ x_0 \in X $, the Picard sequence $ x_n = T^n x_0 $ converges to $ x^* $.
\end{theorem}

\begin{proof}
Let $ x_0 \in X $ be arbitrary. Define $ x_n = T^n x_0 $. Let $ a_k = d(x_k, x_{k+1}) = d(T^k x_0, T^{k+1} x_0) $.

Apply the PA-condition~(\ref{pa1}) to $ x = x_0 $, $ y = T x_0 $. Then:
$ d(T^k x, T^k y) = d(T^k x_0, T^{k+1} x_0) = a_k $, and
$ d(T^{k+1}x, T^{k+1}y) = d(T^{k+1} x_0, T^{k+2} x_0) = a_{k+1} $.

The PA-sum condition gives:
\begin{equation}
\sum_{k=0}^{n-1} a_{k+1} \leq \alpha \sum_{k=0}^{n-1} a_k.
\end{equation}
The left-hand side is $ \sum_{k=1}^{n} a_k $, so:
\begin{equation}
\sum_{k=1}^{n} a_k \leq \alpha \sum_{k=0}^{n-1} a_k.
\end{equation}
Let $ S_n = \sum_{k=0}^{n-1} a_k $. Then:
\begin{equation}
S_n - a_0 + a_n \leq \alpha S_n \implies (1 - \alpha) S_n \leq a_0 - a_n \leq a_0.
\end{equation}
Thus,
\begin{equation}
S_n \leq \frac{a_0}{1 - \alpha} \quad \text{for all } n \ge N.
\end{equation}
Now, define $ C = \max\left\{ S_1, S_2, \dots, S_{N-1}, \frac{a_0}{1 - \alpha} \right\} $. Then
\begin{equation}
S_n \leq C \quad \text{for all } n \in \mathbb{N}.
\end{equation}
Since $ a_k \geq 0 $, $ S_n $ is non-decreasing and bounded, so
\begin{equation}
\sum_{k=0}^{\infty} a_k = \lim_{n \to \infty} S_n < \infty.
\end{equation}
Hence $ \{x_n\} $ is Cauchy (since $ d(x_n, x_{n+1}) = a_n $, summable), so converges to some $ x^* \in X $ by completeness.



Since $ T $ is continuous, $ T x_n \to T x^* $. But $ T x_n = x_{n+1} \to x^* $, so $ T x^* = x^* $.

For uniqueness, suppose $ x^*, y^* $ are fixed points. Then $ d(T^k x^*, T^k y^*) = d(x^*, y^*) $ for all $ k $. For $n \ge N$, apply PA-condition~(\ref{pa1}):
\begin{equation}
\frac{1}{n} \sum_{k=0}^{n-1} d(T^{k+1}x^*, T^{k+1}y^*) = d(x^*, y^*) \leq \alpha \cdot \frac{1}{n} \sum_{k=0}^{n-1} d(T^k x^*, T^k y^*) = \alpha d(x^*, y^*).
\end{equation}
Since $ \alpha < 1 $, $ d(x^*, y^*) = 0 $. Thus, $ x^* = y^* $.
\end{proof}

\begin{remark}
The continuity assumption can be weakened to \textit{orbitally continuity} 
of $ T $~\cite{Ciricorbit,pantorbit}. However, without any continuity, the limit $ x^* $ may not be fixed, as $ T x_n \to T x^* $ cannot be guaranteed.
\end{remark}

\section{Generalization of Banach Contractions}

\begin{proposition}
Every Banach contraction is a PA-contraction.
\end{proposition}

\begin{proof}
Suppose $ d(Tx, Ty) \leq k d(x,y) $ for all $ x,y \in X $, $ k \in (0,1) $. Then for any $ x,y $,
\begin{equation}
d(T^{k+1}x, T^{k+1}y) \leq k d(T^k x, T^k y).
\end{equation}
Summing over $ k = 0 $ to $ n-1 $:
\begin{equation}
\sum_{k=0}^{n-1} d(T^{k+1}x, T^{k+1}y) \leq k \sum_{k=0}^{n-1} d(T^k x, T^k y).
\end{equation}
Thus, $ T $ is a PA-contraction with $ \alpha = k $, $N=1$.
\end{proof}

\begin{example}[PA-contraction not Banach]
Let $ X = \{0, 1, 2\} $ with the discrete metric 
\begin{equation}
d(x,y) = \begin{cases} 0 & x=y \\ 1 & x \ne y \end{cases}.
\end{equation} 
Define:
\begin{equation}
T(0) = 1, \quad T(1) = 2, \quad T(2) = 2.
\end{equation}
$T$ is not a Banach contraction since:
\begin{align*}
d(T0,T1) &= d(1,2) = 1 = d(0,1) \\
d(T0,T2) &= d(1,2) = 1 = d(0,2)
\end{align*}
so no $k < 1$ satisfies $d(Tx,Ty) \leq k d(x,y)$ for all $x,y$.

However, $T$ is a PA-contraction with $\alpha = 1/2$, $N = 2$. For any $x,y \in X$ and $n \geq 2$:
\begin{equation}
\sum_{k=0}^{n-1} d(T^{k+1}x, T^{k+1}y) \leq \frac{1}{2} \sum_{k=0}^{n-1} d(T^k x, T^k y).
\end{equation}

\textbf{Verification for pair (0,1) at n=2:}
\begin{align*}
\text{LHS} &= \sum_{k=0}^{1} d(T^{k+1}0, T^{k+1}1) \\
&= d(T^10, T^11) + d(T^20, T^21) \\
&= d(1,2) + d(2,2) = 1 + 0 = 1 \\
\text{RHS} &= \frac{1}{2} \sum_{k=0}^{1} d(T^k0, T^k1) \\
&= \frac{1}{2} \left[ d(0,1) + d(1,2) \right] \\
&= \frac{1}{2} (1 + 1) = 1
\end{align*}
Since $1 \leq 1$, the condition holds.

\textbf{Verification for pair (0,2) at n=2:}
\begin{align*}
\text{LHS} &= d(T0,T2) + d(T^20,T^22) \\
&= d(1,2) + d(2,2) = 1 + 0 = 1 \\
\text{RHS} &= \frac{1}{2} \left[ d(0,2) + d(1,2) \right] \\
&= \frac{1}{2} (1 + 1) = 1
\end{align*}
$1 \leq 1$ holds.

\textbf{Verification for pair (1,2) at n=2:}
\begin{align*}
\text{LHS} &= d(T1,T2) + d(T^21,T^22) \\
&= d(2,2) + d(2,2) = 0 + 0 = 0 \\
\text{RHS} &= \frac{1}{2} \left[ d(1,2) + d(2,2) \right] \\
&= \frac{1}{2} (1 + 0) = \frac{1}{2}
\end{align*}
$0 \leq 1/2$ holds.

For $n > 2$ and any pair, $T^k x = T^k y = 2$ for $k \geq 2$, so all distances are 0 for $k \geq 2$. The sums are thus identical to the n=2 case.
Hence, $ T $ is a PA-contraction but not a Banach contraction.
\label{ex:panotbanach}
\end{example}



\section{Comparison with F-Contractions}

We recall Wardowski's F-contraction.

\begin{definition}[F-Contraction \cite{Wardowski2012}]
A mapping $ T: X \to X $ is an \textbf{F-contraction} if there exist $ \tau > 0 $ and a function $ F: (0,\infty) \to \mathbb{R} $ satisfying
\begin{enumerate}
\item[(F1)] $ F $ is strictly increasing,
\item[(F2)] $ \lim_{t \to 0^+} F(t) = -\infty $,
\item[(F3)] $ \lim_{t \to 0^+} t^k F(t) = 0 $ for some $ k \in (0,1) $,
\end{enumerate}
such that for all $ x,y \in X $ with $ d(Tx,Ty) > 0 $,
\begin{equation}
\tau + F(d(Tx, Ty)) \leq F(d(x,y)).
\end{equation}
\end{definition}



\begin{example}[F-contraction not PA-contraction]

No known simple example. But since F-contractions are pointwise and PA-contractions are averaged, they are conceptually distinct.
The relationship between F- and PA-contractions is an open problem.
\end{example}

%%%%%%%%%%%%%%%%%%%%%%%%%%%%%%%%%%%%%
%%%%%%%%%%%%%%%%%%%%%%%%%%%%%%%%%%%%%%


\begin{example}[PA-contraction not F-contraction]
Let $ T x = x^2/2 $ on $ [0,1] $. We show that $ T $ is a PA-contraction but not a  F-contraction.

First, we have that
\begin{equation}
T^k x = \frac{x^{2^k}}{2^{2^k - 1}}.
\label{t2}
\end{equation}
For any $ x, y \in [0,1] $, the distance $ d(T^k x, T^k y) $ decays double-exponentially as $ k \to \infty $. 
For $Tx = x^2/2$, explicit computation shows
the uniform bound for $n \ge 5$ and shows that, for $ n \geq 5 $, the ratio
\begin{equation}
\frac{ \sum_{k=0}^{n-1} d(T^{k+1}x, T^{k+1}y) }{ \sum_{k=0}^{n-1} d(T^k x, T^k y) }
\end{equation}
is uniformly bounded by $ \alpha = 0.4 $ across all $ x, y \in [0,1] $. Thus, $ T $ is a PA-contraction with $ N = 5 $.


Finally, as $ x,y \to 1^- $, $ \frac{d(Tx,Ty)}{d(x,y)} \to 1 $, so $ T $ is not a Banach contraction. For any $ F $ satisfying Wardowski's conditions (F1)--(F3), the inequality $ \tau + F(d(Tx,Ty)) \leq F(d(x,y)) $ fails for $ \tau > 0 $ in this limit. Hence, $ T $ is not an F-contraction.
\end{example}

%%%%%%%%%%%%%%%%%%%%%%%%%%%%%%%%%%%%%
%%%%%%%%%%%%%%%%%%%%%%%%%%%%%%%%%%%%%%


%%%%%%%%%%%%%%%%%%%%%%%%%%%%%%%%%%%%%%%%%%%%%
%%%%%%%%%%%%%%%%%%%%%%%%%%%%%%%%%%%%%%%%%%%%%%%
%%%%%%%%%%%%%%%%%%%%%%%%%%%%%%%%%%%%%%%%%%%%%%%


%%%%%%%%%%%%%%%%%%%%%%%%%%%%%%%%%%%%%%%%%%%%%%%%%%%%%
%%%%%%%%%%%%%%%%%%%%%%%%%%%%%%%%%%%%%%%%%%%%%%%%%%%%%
%%%%%%%%%%%%%%%%%%%%%%%%%%%%%%%%%%%%%%%%%%%%%%%%%%%%%


\section{Relation with Other Classical Contractions}

We now clarify the relationship between PA-contractions and other classical generalizations.

\begin{remark}[Kannan Contractions]
A Kannan contraction satisfies $ d(Tx,Ty) \leq k[d(x,Tx) + d(y,Ty)] $, $ k < 1/2 $. While Kannan maps are asymptotically regular and generate summable sequences $ \sum d(T^n x, T^{n+1}x) < \infty $, they do not necessarily satisfy the uniform PA-condition for all $ n $. The ratio of averages may exceed $ \alpha < 1 $ for small $ n $. 

Next, we show that $ T $ defined in~(\ref{t2}), a PA-contraction, is not a Kannan contraction. Suppose, for contradiction, that
\begin{equation}
d(Tx,Ty) \leq k \left[ d(x,Tx) + d(y,Ty) \right], \quad k < \frac{1}{2}.
\end{equation}
Take $ x = 1 $, $ y = 0 $. Then:
 $ d(Tx,Ty) = |1/2 - 0| = 1/2 $,
 $ d(x,Tx) = |1 - 1/2| = 1/2 $,
 $ d(y,Ty) = |0 - 0| = 0 $.
 Right-hand side = $ k(1/2 + 0) = k/2 $.

Then $ 1/2 \leq k/2 \Rightarrow k \geq 1 $, contradicting $ k < 1/2 $. So $ T $ is not a Kannan contraction.


Thus, \textbf{Kannan contractions are not generally PA-contractions}, 
\textbf{PA-contractions are not generally Kannan contractions}, 
so the classes are incomparable.
\label{rem:kannan}
\end{remark}



\begin{remark}[Chatterjea not PA-contraction]
A mapping $ T: X \to X $ on a complete metric space $ (X,d) $ is a Chatterjea contraction~\cite{Chatterjea1972} if there exists $ k \in (0, \frac{1}{2}) $ such that for all $ x, y \in X $,
\begin{equation}
d(Tx, Ty) \leq k \left[ d(x, Ty) + d(y, Tx) \right].
\label{chatterjea}
\end{equation}
Let $ X = \mathbb{N} \cup \{\infty\} $, $ d(m,n) = |1/m - 1/n| $, $ d(n,\infty) 
= 1/n $, $ Tn = n+1 $, $ T\infty = \infty $. 
The PA-Contraction is defined in~(\ref{pa1}). We will show that $T$ is Chatterjea but not PA.


$ T $ is a Chatterjea Contraction.

Let $ m, n \in \mathbb{N} $, $ m \ne n $. We compute both sides.
Left-hand side
\begin{equation}
d(Tm, Tn) = \left| \frac{1}{m+1} - \frac{1}{n+1} \right| \; .
\end{equation}
Right-hand side
\begin{equation}
d(m, Tn) + d(n, Tm) = \left| \frac{1}{m} - \frac{1}{n+1} \right| + \left| \frac{1}{n} - \frac{1}{m+1} \right| \; .
\end{equation}
We claim that
\begin{equation}
\left| \frac{1}{m+1} - \frac{1}{n+1} \right| \leq k \left( \left| \frac{1}{m} - \frac{1}{n+1} \right| + \left| \frac{1}{n} - \frac{1}{m+1} \right| \right)
\quad \text{for some } k < \frac{1}{2} \;.
\end{equation}
Consider asymptotic behavior as $ m, n \to \infty $,
let $ m = n $. Then
$ d(Tn, Tn) = 0 $, so inequality holds.

Let $ m = n+1 $. Then
$ d(Tm, Tn) = \left| \frac{1}{n+2} - \frac{1}{n+1} \right| = \frac{1}{(n+1)(n+2)} $,
$ d(m, Tn) = \left| \frac{1}{n+1} - \frac{1}{n+1} \right| = 0 $,
$ d(n, Tm) = \left| \frac{1}{n} - \frac{1}{n+2} \right| = \frac{2}{n(n+2)} $.
So RHS = $ 0 + \frac{2}{n(n+2)} $, LHS = $ \frac{1}{(n+1)(n+2)} $.

Then
\begin{equation}
\frac{\text{LHS}}{\text{RHS}} = \frac{ \frac{1}{(n+1)(n+2)} }{ \frac{2}{n(n+2)} } = \frac{n}{2(n+1)} \to \frac{1}{2}
\quad \text{as } n \to \infty \; .
\end{equation}
So for large $ n $, $ d(Tm,Tn) \leq \left( \frac{1}{2} - \varepsilon_n \right) [d(m,Tn) + d(n,Tm)] $, with $ \varepsilon_n \to 0 $.
Thus, for any $ k > \frac{1}{2} $, it fails, but for $ k = \frac{1}{2} - \delta $, we can find $ N $ such that for $ n > N $, it holds.

For small $ n $, since $ X $ is discrete, the ratio is bounded, so we can take $ k = \frac{1}{2} - \delta $ for small $ \delta > 0 $, and verify finitely many cases.

Hence, $ T $ satisfies the Chatterjea condition with some $ k < \frac{1}{2} $.
So $ T $ is a Chatterjea contraction.

%%%%%%%%%%%%%%%%%%%%%%%%%%%%%%%%%%%%%%%%%
To prove that $T$ is not a PA-contraction
we will have to show that there is no $ \alpha \in (0,1) $ such that for all $ N \in \mathbb{N} $ and all $ m, n \in \mathbb{N} $,
\begin{equation}
  \frac{1}{n} \sum_{k=0}^{n-1} d(T^{k+1}n, T^{k+1}m) \leq \alpha \cdot \frac{1}{n} \sum_{k=0}^{n-1} d(T^k n, T^k m).
\end{equation}
We will show that the ratio of averages approaches 1 as $ n \to \infty $, so no uniform $ \alpha < 1 $ works.

We have
\begin{equation}
T^k n = n + k, \quad T^k m = m + k
\Rightarrow
d(T^k n, T^k m) = \left| \frac{1}{n+k} - \frac{1}{m+k} \right| \; .
\end{equation}
Assume without loss of generality that $ n < m $. Then $ \frac{1}{n+k} > \frac{1}{m+k} $, so
\begin{equation}
d(T^k n, T^k m) = \frac{1}{n+k} - \frac{1}{m+k} \; .
\end{equation}
Define the average distance over $ k = 0 $ to $ n-1 $
\begin{equation}
A_n(n,m) := \frac{1}{n} \sum_{k=0}^{n-1} d(T^k n, T^k m) = \frac{1}{n} \sum_{k=0}^{n-1} \left( \frac{1}{n+k} - \frac{1}{m+k} \right) \; ,
\end{equation}
and similarly, the next average shifted by one iterate
\begin{equation}
A_n^{(1)}(n,m) := \frac{1}{n} \sum_{k=0}^{n-1} d(T^{k+1}n, T^{k+1}m) = \frac{1}{n} \sum_{k=0}^{n-1} \left( \frac{1}{n+k+1} - \frac{1}{m+k+1} \right)
= \frac{1}{n} \sum_{k=1}^{n} \left( \frac{1}{n+k} - \frac{1}{m+k} \right) \; .
\end{equation}
So
\begin{equation}
A_n^{(1)}(n,m) = \frac{1}{n} \sum_{k=1}^{n} \left( \frac{1}{n+k} - \frac{1}{m+k} \right)
= A_n(n,m) - \frac{1}{n} \left( \frac{1}{n+0} - \frac{1}{m+0} \right) + \frac{1}{n} \left( \frac{1}{n+n} - \frac{1}{m+n} \right) \; 
\end{equation}
that is
\begin{equation}
A_n^{(1)} = A_n - \frac{1}{n} \left( \frac{1}{n} - \frac{1}{m} \right) + \frac{1}{n} \left( \frac{1}{2n} - \frac{1}{m+n} \right) \; .
\end{equation}
We are interested in whether
\begin{equation}
\frac{A_n^{(1)}}{A_n} \leq \alpha < 1 \quad \text{uniformly in } n
\end{equation}
for some  $\alpha < 1$  and all  $n, m$.


To analyze asymptotic behavior, fix $ m = n + 1 $. Then
\begin{equation}
d(T^k n, T^k m) = \left| \frac{1}{n+k} - \frac{1}{n+1+k} \right| = \frac{1}{n+k} - \frac{1}{n+k+1} = \frac{1}{(n+k)(n+k+1)} \; .
\end{equation}

So
\begin{equation}
A_n = \frac{1}{n} \sum_{k=0}^{n-1} \frac{1}{(n+k)(n+k+1)} \; ,
\end{equation}
observe that $ n+k $ runs from $ n $ to $ 2n-1 $. Thus,
\begin{equation}
A_n = \frac{1}{n} \sum_{j=n}^{2n-1} \frac{1}{j(j+1)} \; .
\end{equation}
Now use
\begin{equation}
\frac{1}{j(j+1)} = \frac{1}{j} - \frac{1}{j+1}
\Rightarrow \sum_{j=n}^{2n-1} \left( \frac{1}{j} - \frac{1}{j+1} \right) = \frac{1}{n} - \frac{1}{2n} = \frac{1}{2n} \;,
\end{equation}
so
\begin{equation}
A_n = \frac{1}{n} \cdot \frac{1}{2n} = \frac{1}{2n^2} \; .
\end{equation}
Compute $ A_n^{(1)} = \frac{1}{n} \sum_{k=0}^{n-1} d(T^{k+1}n, T^{k+1}m) = \frac{1}{n} \sum_{k=1}^{n} \frac{1}{(n+k)(n+k+1)} $, where
$ j = n+k $ runs from $ n+1 $ to $ 2n $, 
\begin{equation}
\sum_{j=n+1}^{2n} \frac{1}{j(j+1)} = \sum_{j=n+1}^{2n} \left( \frac{1}{j} - \frac{1}{j+1} \right) = \frac{1}{n+1} - \frac{1}{2n+1} \; .
\end{equation}
Thus
\begin{equation}
A_n^{(1)} = \frac{1}{n} \left( \frac{1}{n+1} - \frac{1}{2n+1} \right)
= \frac{1}{n} \left( \frac{(2n+1) - (n+1)}{(n+1)(2n+1)} \right)
= \frac{1}{n} \cdot \frac{n}{(n+1)(2n+1)}
= \frac{1}{(n+1)(2n+1)} \; .
\end{equation}
Now calculate the ratio
\begin{equation}
\frac{A_n^{(1)}}{A_n} = \frac{ \frac{1}{(n+1)(2n+1)} }{ \frac{1}{2n^2} } = \frac{2n^2}{(n+1)(2n+1)} \; ,
\end{equation}
and simplify
\begin{equation}
= \frac{2n^2}{2n^2 + n + 2n + 1} = \frac{2n^2}{2n^2 + 3n + 1}
= \frac{2}{2 + \frac{3}{n} + \frac{1}{n^2}} \to 1 \quad \text{as } n \to \infty \; ,
\end{equation}
so we conclude that
\begin{equation}
\lim_{n \to \infty} \frac{A_n^{(1)}}{A_n} = 1 \; .
\end{equation}
Therefore, for any $ \alpha < 1 $, there exists $ n $ large enough such that:
\begin{equation}
A_n^{(1)} > \alpha A_n \; ,
\end{equation}
hence, the PA-condition~(\ref{pa1})
\begin{equation}
\frac{1}{n} \sum_{k=0}^{n-1} d(T^{k+1}n, T^{k+1}m) \leq \alpha \cdot \frac{1}{n} \sum_{k=0}^{n-1} d(T^k n, T^k m)
\end{equation}
cannot hold uniformly for all $ n \in \mathbb{N} $ with a fixed $ \alpha < 1 $.
Thus, $T$ is not a PA-contraction.


 
\label{rem:chatterjeanotpa}
\end{remark}

\begin{remark}[\'Ciri\'c not PA-contraction]
A mapping $ T: X \to X $ on a complete metric space $ (X,d) $ is a \'Ciri\'c contraction if there exists $ k \in (0,1) $ such that for all $ x, y \in X $,
\begin{equation}
d(Tx, Ty) \leq k \cdot \max\left\{
d(x,y),\;
d(x,Tx),\;
d(y,Ty),\;
\frac{d(x,Ty) + d(y,Tx)}{2}
\right\}.
\label{ciric}
\end{equation}
The PA-Contraction is defined in~(\ref{pa1}). We will use the same definitions of Example~\ref{rem:chatterjeanotpa} and show that $T$ is \'Ciri\'c but not PA.

$ T $  is a \'Ciri\'c contraction.

Let $ m, n \in \mathbb{N} $, $ m \ne n $. Without loss of generality, assume $ m < n $.

We compute
$ d(Tm, Tn) = \left| \frac{1}{m+1} - \frac{1}{n+1} \right| $,
$ d(m,n) = \left| \frac{1}{m} - \frac{1}{n} \right| $,
$ d(m, Tm) = \left| \frac{1}{m} - \frac{1}{m+1} \right| = \frac{1}{m(m+1)} $,
$ d(n, Tn) = \frac{1}{n(n+1)} $,
$ d(m, Tn) = \left| \frac{1}{m} - \frac{1}{n+1} \right| $,
$ d(n, Tm) = \left| \frac{1}{n} - \frac{1}{m+1} \right| $,
$ \frac{d(m,Tn) + d(n,Tm)}{2} $.

We claim that $ d(Tm, Tn) \leq k \cdot \max\{ \dots \} $ for some $ k < 1 $.


Consider the full max of~(\ref{ciric}).

For large $ m, n $, $ d(m,Tm) = \frac{1}{m(m+1)} \approx \frac{1}{m^2} $, very small,
and $ d(m,n) \approx \frac{|m-n|}{m^2} $ if $ m \approx n $, also small.
And $ d(Tm, Tn) \approx \frac{|m-n|}{(m+1)^2} $.

Now consider $ \frac{d(m,Tn) + d(n,Tm)}{2} $. If $ m < n $, $ Tn = n+1 < n $, so $ \frac{1}{m} > \frac{1}{n+1} $, so
\begin{equation}
d(m,Tn) = \frac{1}{m} - \frac{1}{n+1}, \quad
d(n,Tm) = \frac{1}{n} - \frac{1}{m+1} \quad \text{(if } m+1 < n\text{)} \; .
\end{equation}
So
\begin{gather}
\frac{d(m,Tn) + d(n,Tm)}{2} = \frac{1}{2} \left( \frac{1}{m} - \frac{1}{n+1} + \frac{1}{n} - \frac{1}{m+1} \right) 
= \frac{1}{2} \left( \left( \frac{1}{m} - \frac{1}{m+1} \right) + \left( \frac{1}{n} - \frac{1}{n+1} \right) \right) = \nonumber \\
 \frac{1}{2} \left( \frac{1}{m(m+1)} + \frac{1}{n(n+1)} \right) \; ,
\end{gather}
and this is small.

$ d(m,n) = \frac{1}{m} - \frac{1}{n} = \frac{n - m}{mn} $,
and $ d(Tm,Tn) = \frac{n - m}{(m+1)(n+1)} $.
So
\begin{equation}
\frac{d(Tm,Tn)}{d(m,n)} = \frac{n - m}{(m+1)(n+1)} \cdot \frac{mn}{n - m} = \frac{mn}{(m+1)(n+1)} < 1
\end{equation}


Let $ k = \sup \frac{mn}{(m+1)(n+1)} $. Since this $ < 1 $ for all $ m,n $, and $ \to 1 $ as $ m,n \to \infty $, we can pick any $ k \in (0,1) $, say $ k = 0.99 $, and for all $ m,n $ such that $ \frac{mn}{(m+1)(n+1)} \leq k $, it holds.

For small $ m,n $, verify finitely many cases: since all values are bounded away from equality, we can always find a $ k < 1 $ such that:
\begin{equation}
d(Tm,Tn) \leq k \cdot d(m,n) \leq k \cdot M(m,n) \; .
\end{equation}

Thus, $ T $ is a \'Ciri\'c contraction with $ k < 1 $.


To prove that $ T $ is not a PA-Contraction, fix  $ m, n \in \mathbb{N} $,
 $ m = n - 1 $. The proof goes exactly as in Remark~\ref{rem:chatterjeanotpa}.


\label{rem:ciricnotpa}
\end{remark}

Whether PA-contractions imply Chatterjea or \'Ciri\'c contractions remains open. Conversely, Remarks~\ref{rem:chatterjeanotpa} and~\ref{rem:ciricnotpa} show they do not.

\section{Comparison Table}

The following table summarizes key properties of major contraction types.

\begin{center}
\scalebox{0.8}{
\begin{tabular}{lcccccc}
\toprule
\textbf{Contraction Type} & \textbf{Banach} & \textbf{Kannan} & \textbf{Chatterjea} & \textbf{\'Ciri\'c} & \textbf{F} & \textbf{PA} \\
\midrule
Pointwise inequality & Yes & Yes & Yes & Yes & Yes & No (averaged) \\
Requires continuity & No & No & No & No & No & Yes (for fixed point) \\
Generalizes Banach & No & Yes & Yes & Yes & Yes & Yes \\
Summable $ d(T^n x, T^{n+1}x) $ & Yes & Yes & Yes & Yes & Sometimes & Yes \\
Based on orbits/averages & No & No & No & No & No & \textbf{Yes} \\
Implies PA & Yes & No & No & No & Open & — \\
PA implies & No & No & Open & Open & No & — \\
\bottomrule
\end{tabular}}
\end{center}

\begin{remark}
While we have shown that PA-contractions are not implied by classical types, the converse — whether PA-contractions imply other classes — is only partially resolved. The mapping $ T x = x^2/2 $ shows that PA-contractions do not imply Kannan or F-contractions. Whether every F-contraction is a PA-contraction, and if every 
PA-contraction is a Chatterjea or \'Ciri\'c contraction remains open.
\end{remark}

\section{Conclusion}

We introduced \textbf{PA-contractions}, a new class based on averaged contraction over orbits. This condition strictly generalizes Banach contractions,
is independent of F-contractions, Kannan, Chatterjea, and \'Ciri\'c types,
ensures unique fixed points under continuity. Future investigation could go
in the directions of solving open problems, multivalued mappings and common fixed points.

PA-contractions offer a new lens for analyzing asymptotic behavior in fixed point theory.
Our work continues the long tradition of generalizing the Banach principle, as studied by Rhoades \cite{Rhoades1977} and extended by Wardowski \cite{Wardowski2012}, Proinov \cite{Proinov2020}, Ran and 
Reurings~\cite{Ran2004}, Alam and Imdad~\cite{Alam2017}, and others.


%\bibliographystyle{plain}
\begin{thebibliography}{33}



\bibitem{Agarwal2011}
Agarwal, R. P., Meehan, M., \& O'Regan, D. (2001). \emph{Fixed Point Theory and Applications}. Cambridge University Press.
\url{http://dx.doi.org/10.1017/cbo9780511543005}

\bibitem{Alam2017}
Alam, A., \& Imdad, M. (2015). \emph{Relation-theoretic contraction principle.} Journal of Fixed Point Theory and Applications, 19(2), 1451–1474.
\url{http://dx.doi.org/10.1007/s11784-015-0247-y}

\bibitem{Banach1922} Banach, S. (1922). \emph{Sur les opérations dans les ensembles abstraits et leur application aux équations intégrales.} Fund. Math., 3, 133–181.


\bibitem{Ciric1971} \'Ciri\'c, L. B. (1971). \emph{Generalized contractions and fixed-point theorems}. Publ. Inst. Math., 12(26), 19–26.
\bibitem{Ciricorbit} \'Ciri\'c, L. B. (1971). \emph{On contraction type mappings}, Math. Balkanica 1, 52-57.
\bibitem{Ciric1974}  \'Ciri\'c, L. B., \emph{A generalization of Banach’s contraction principle}, Proc. Amer. Math. Soc. 45 (1974), 267–273.
\bibitem{Ciric1979} \'Ciri\'c, L. B., \emph{On mappings with contractive iteration},  Publ. Inst. Math.(N.S) 46 (40) (1979), 79–82
\bibitem{Chatterjea1972} Chatterjea, S.K. 1972. \emph{Fixed-point theorems.} 
Comptes Rendus de l’Academie bulgare des Sciences, 25(6), pp. 727–730.

\bibitem{fabianopathintegral}
 Fabiano, N. (2022). \emph{Path integral in quantum field theories}, Vojnotehni\v{c}ki glasnik/Military Technical Courier, 70(4),  pp. 993-1016. 
 \url{http://dx.doi.org/10.5937/vojtehg70-35882}
%doi:10.5937/vojtehg70-35882.
 

\bibitem{fabiano2022}
Fabiano N., Kadelburg Z., Mirkov N., \v{S}e\v{s}um \v{C}avi\'c V., Radenovi\'c S. On \emph{F-Contractions: A Survey}. Contemp. Math. [Internet]. 2022 Aug. 9 [cited 2025 Jul. 24];3(3):327-42. 
\url{https://doi.org/10.37256/cm.3320221517}
%Available from: https://ojs.wiserpub.com/index.php/CM/article/view/1517

\bibitem{Hussain2012}
Hussain, N., Kutbi, M. A., \& Salimi, P. (2012). \emph{Fixed point theorems for 
$\alpha$-$\psi$-contractive type mappings}. Fixed Point Theory and Applications, 2014(1), 124.
\url{https://doi.org/10.1155/2014/280817}

\bibitem{Kannan1968} Kannan, R. (1968). \emph{Some results on fixed points}. Bull. Calcutta Math. Soc., 60, 71–76.


\bibitem{pantorbit} Pant, A. and Pant, R.P., (2017). \emph{Orbital continuity and fixed points}. Filomat, 31(11), pp.3495-3499.
\url{http://dx.doi.org/10.2298/fil1711495p}

\bibitem{Piri2014}
Piri, H., \& Kumam, P. (2014). \emph{Some fixed point theorems concerning 
F-contraction in complete metric spaces}. Fixed Point Theory and Applications, 2014(1), 210.
\url{http://dx.doi.org/10.1186/1687-1812-2014-210}


\bibitem{Proinov2020}
Proinov, P. D. (2020). \emph{Fixed point theorems for generalized contractive mappings in metric spaces}. Journal of Fixed Point Theory and Applications, 22(1), 21.
\url{http://dx.doi.org/10.1007/s11784-020-0756-1}


\bibitem{Ran2004}
Ran, A. C. M., \& Reurings, M. C. B. (2003). \emph{A fixed point theorem in partially ordered sets and some applications to matrix equations}. Proceedings of the American Mathematical Society, 132(5), 1435–1443.
\url{http://dx.doi.org/10.1090/s0002-9939-03-07220-4}

\bibitem{Rhoades1977}
Rhoades, B. E. (1977). \emph{A comparison of various definitions of contractive mappings}. Transactions of the American Mathematical Society, 226, 257–290.
\url{http://dx.doi.org/10.2307/1997954}

\bibitem{Secelean2013}
Secelean, N. A. (2013). \emph{Iterated function systems consisting of 
F-contractions}. Fixed Point Theory and Applications, 2013(1), 277.
\url{http://dx.doi.org/10.1186/1687-1812-2013-277}


\bibitem{Wardowski2012} Wardowski, D. (2012). \emph{Fixed points of a new type of contractive mappings in complete metric spaces.} Fixed Point Theory Appl., 2012(1), 94.
\url{http://dx.doi.org/10.1186/1687-1812-2012-94}

\end{thebibliography}

\end{document}
